% v2.7.0 figure UID: FIG-P654-01
% Image-guided reverse redraw: preserve the 8-node/7-path dependency graph and separate all lower-row labels.
\tikzset{slfig-FIG-P654-01/.style={font=\fontsize{10.1pt}{12.2pt}\selectfont,every path/.append style={line cap=round,line join=round}}}
\pgfkeysifdefined{/tikz/alt}{}{\tikzset{alt/.code={}}}
\begingroup
% One 11.6pt BASE_MATH role: IBM Plex math letters, a Libertinus U+002B
% binary operator, and a genuine Mezenets U+004E uppercase N.  No
% per-glyph font-size overrides or semantic substitutions.
\font\slfigPBaseMath="IBMPlexMath-Regular.otf" at 11.6pt
\font\slfigPBinaryMath="LibertinusMath-Regular.otf" at 11.6pt
\font\slfigPTotalMath="MezenetsUnicode.otf" at 11.6pt
\def\slfigPAlpha{\mathord{\hbox{\slfigPBaseMath\char"1D6FC}}}
\def\slfigPCountN{\mathord{\hbox{\slfigPBaseMath\char"1D45B}}}
\def\slfigPPlus{\mathbin{\hbox{\slfigPBinaryMath\char"002B}}}
\def\slfigPTotalN{\mathord{\hbox{\slfigPTotalMath\char"004E}}}
\begin{figure}[htbp]
\centering
\begin{tikzpicture}[slfig-FIG-P654-01,
  >=Stealth,
  every node/.style={font=\fontsize{10.1pt}{12.2pt}\selectfont,align=center},
  core/.style={draw=SLMidBlue,fill=SLSoftBlue,rounded corners=2pt,
    minimum height=13mm,text width=31mm,inner xsep=4pt,inner ysep=3.8pt},
  aux/.style={draw=SLRuleGray,fill=SLSoftGray,rounded corners=2pt,
    minimum height=12mm,text width=31mm,inner xsep=4pt,inner ysep=3.8pt},
  result/.style={draw=SLDeepBlue,fill=white,rounded corners=2pt,
    minimum height=18mm,text width=37mm,inner xsep=3pt,inner ysep=4pt},
  arr/.style={->,draw=SLMidBlue,line width=.95pt},
  interp/.style={draw=SLGray,line width=.62pt},
  application/.style={->,draw=SLGray,line width=.70pt,dash pattern=on 2.2pt off 1.5pt},
  >={Stealth[length=2.4mm,width=1.7mm]},
  every node/.append style={line width=0.85pt},
  alt={对象—关系—结论：类别计数与Gamma—Beta归一化两条前置线索汇入多项似然和Dirichlet先验，实线有向主链先以参数相加得到共轭后验，再由后验均值得到给定计数与超参数的类别预测概率；无箭头细线连接单纯形几何、均值浓度和对数矩等解释对象，虚线有向边仅把预测连到下游主题模型应用}]
\node[aux,text width=28mm] (trial) at (-5.499,1.15) {类别计数\\{\fontsize{11.6pt}{13.8pt}\selectfont$\slfigPCountN$}};
\node[aux,text width=28mm] (gamma) at (-5.499,-1.15) {Gamma 与贝塔\\规范因子};
\node[core] (families) at (-1.88,0)
  {多项分布\\与狄利克雷先验};
\node[result] (posterior) at (2.15,0)
  {共轭狄利克雷后验\\[5pt]
   {\fontsize{11.6pt}{13.8pt}\selectfont
    $\displaystyle\text{参数 }\slfigPAlpha\slfigPPlus\slfigPCountN$}};
\node[result,text width=43mm,minimum height=30mm] (predictive) at (6.65,0)
  {后验预测概率\\
   新增观测取指定类别\\[5pt]
   {\fontsize{11.6pt}{13.8pt}\selectfont
    $\displaystyle\dfrac{\slfigPAlpha_i\slfigPPlus\slfigPCountN_i}{\slfigPAlpha_0\slfigPPlus\slfigPTotalN}$}};
\node[aux,text width=23.5mm,inner xsep=2.5pt] (simplex) at (-1.41,-2.12) {单纯形几何};
\node[aux,text width=23.5mm,inner xsep=2.5pt] (mom) at (2.0,-2.12) {均值与浓度\\及对数矩};
\node[aux,text width=23.5mm,inner xsep=2.5pt] (lda) at (6.35,-2.75)
  {后续主题模型\\伪计数与平滑};

\draw[arr] (trial) -- (families);
\draw[arr] (gamma) -- (families);
\draw[arr] (families) -- (posterior);
\draw[arr] (posterior) -- (predictive);
\draw[interp] (families)--(simplex);
\draw[interp] (posterior)--(mom);
\draw[application] (predictive)--node[pos=.50,right=7pt,inner sep=0pt,
  font=\fontsize{10.1pt}{12.2pt}\selectfont,text=SLTextGray]{应用}(lda);
\end{tikzpicture}
\captionsetup{width=.94\linewidth}
\caption{Dirichlet--多项共轭更新与单步后验预测主链}
\label{fig:V5-C05-dependency-graph}
\end{figure}
\endgroup
